\documentclass[10.5pt]{letter}

\usepackage[a4paper,margin=0.8in]{geometry}
\usepackage{hyperref}

\hypersetup{
    colorlinks=true,
    urlcolor=blue
}

\setlength{\parskip}{0.4em}
\setlength{\parindent}{0pt} 

\signature{
  \textbf{Mohd Anas} \\
  B.Tech Information Technology, IIIT Vadodara ('27)
}

\begin{document}

\begin{letter}{
  From: \\
  Mohd  Anas \\
  anas.ahamad955@gmail.com $|$ +91 8081233871 $|$ \href{https://linkedin.com/in/anas-um}{LinkedIn} $|$ \href{https://github.com/Anas-github-acc}{Github} \\
  To: \\
  Hiring Team \\
  Stripe
}

\opening{Dear Hiring Team, }
I am writing to apply for the Software Engineer Intern position at Stripe. I am currently in my third year of pursuing a B.Tech degree in Information Technology at IIIT Vadodara and will be graduating in 2027. Throughout the past year, I have got ample opportunity to work on backend services, distributed systems, and production software. It is specifically due to this reason that I am eager to intern at Stripe wherein interns get to work on production software and services which customers use.

While interning at Aarogya Kaya, I got a chance to work extensively with Go programming language. I got to build services integrating with the APIs. Moreover, I also got to work on CQRS, PostgreSQL, Redis, and in-memory caching. One of the challenges I addressed was performance. By optimizing the approach to caching and fetching data from the database, I successfully managed to bring down average read time from 420 ms to 38 ms while maintaining p95 latency less than 250 ms.

Outside of my internship I enjoy building systems that force me to think about what happens when things fail. I built a distributed key-value store with replication, gossip-based membership, failure detection, synchronization and anti-entropy repair across a three-node cluster. I deployed it on AWS EKS using Kubernetes StatefulSets and built automated tests for replication, synchronization and node-failure scenarios. Working on it gave me a deeper understanding of distributed systems because I had to think beyond the happy path—how nodes recover, how replicas stay consistent and how the system behaves when part of it becomes unavailable.

I have also contributed to open-source projects, including Next.js. Have had 22 pull requests merged across different projects. Working in codebases maintained by engineers taught me how to understand unfamiliar systems respond to review feedback and make changes that fit existing engineering conventions rather than simply implementing something in isolation. That experience is one of the reasons Stripe’s emphasis on code reviews, collaboration and learning from engineers stood out to me.

What excites me most about Stripe is the engineering environment behind the product. Payments may look simple from the outside. The systems supporting them need to be reliable, scalable, understandable and correct. Those are exactly the kinds of engineering problems I want to spend time working on. I would be excited to contribute what I already know about backend and distributed systems while learning from engineers who build infrastructure operating at Stripe’s scale.

I would appreciate the opportunity to contribute to Stripe as an intern and to ship software that has an impact, on its users.

Thank you for your time and consideration.
\closing{
  Sincerely,
}


\end{letter}

\end{document}
